\section{Discussion}

\subsection{Scalar T/I/F Is Necessary But Insufficient}

The original Smarandache and Leyva-Vázquez thesis is correct:
independent T/I/F dimensions reveal hyper-truth that probabilistic
constraints suppress. Our cross-vendor parallel experiment extends this
claim across five vendors and reports higher S1 hyper-truth frequency
under our independently constructed prompts.

However, scalar T/I/F has a ceiling. The Absorption problem shows that
some models cannot express the distinction between different types of
uncertainty using three numbers alone. This is not a model deficiency,
it is a representational limitation. The model carries the distinction
internally but lacks the output dimensions to express it.

\subsection{Declared Losses Are the Missing Dimension}

Adding structured loss declarations to the evaluation output recovers
the collapsed distinctions. The loss channel carries semantic
information that the scalar channel drops: \emph{why} the model is
uncertain, \emph{what} it cannot evaluate, and \emph{how severely} this
limits its assessment.

This suggests a hierarchy of epistemic output fidelity:

\begin{enumerate}
\def\labelenumi{\arabic{enumi}.}
\tightlist
\item
  \textbf{Probabilistic} (T+I+F=1): Collapses all uncertainty to a
  single dimension. Cannot express conflict. Cannot express hyper-truth.
\item
  \textbf{Scalar Neutrosophic} (independent T,I,F): Expresses conflict
  and hyper-truth. Cannot distinguish types of uncertainty for
  Absorption models.
\item
  \textbf{Tensor Neutrosophic} (T,I,F + declared losses): Expresses
  conflict, hyper-truth, and type-specific uncertainty. Recovers
  distinctions that scalars collapse.
\end{enumerate}

Each level subsumes the previous. The tensor level is the minimum needed
for faithful representation of LLM epistemic state.

\subsection{Implications for Evaluation Architecture}

The finding that models can produce accurate, domain-specific loss
declarations has practical implications for AI evaluation systems. An
attestation system that records only scalar T/I/F values will lose
information that loss declarations preserve. Systems designed for
routing evaluations to appropriate evaluators by selecting which model
or method to use for a given claim, should route on losses rather than
capabilities, because losses carry richer information about what the
evaluator cannot do and why.

\subsection{Limitations}

\begin{enumerate}
\def\labelenumi{\arabic{enumi}.}
\item
  \textbf{Prompt sensitivity}: We do not have the original study's
  prompts, so the S1--S3 comparison is across independently constructed
  prompts. The S2 agreement provides evidence of compatibility but not
  identity.
\item
  \textbf{Sample size}: 5 repetitions per cell provides stability
  measures but is limited for confirmatory hypothesis testing under
  strict independence assumptions. The statistical tests in
  Section~\ref{sec:validation} should be read as exploratory evidence.
  The
  zero-variance findings (e.g., Claude's paradox position) are robust;
  the variable findings (e.g., Qwen's paradox instability) would benefit
  from 30+ repetitions.
\item
  \textbf{Loss vocabulary metric}: Jaccard similarity on word-level
  tokens is a crude measure of semantic differentiation. More
  sophisticated NLP measures (sentence embeddings, BERTScore) might
  reveal structure that word-level overlap misses or might show that
  some apparently disjoint vocabularies express similar concepts.
\item
  \textbf{Self-report vs.~internal state}: The declared losses are the
  model's self-report of its limitations, not a direct observation of
  internal state. Models might produce plausible-sounding but generic
  loss declarations without genuine internal distinction. The
  domain-specificity and consistency of the losses across repetitions
  argues against this, but independent validation (e.g., via logprob
  analysis) would strengthen the finding.
\item
  \textbf{Five phenomena, five models}: The stimuli set is small
  (inherited from the original study). A broader range of epistemic
  phenomena would test the generality of both the Absorption problem and
  the tensor solution.
\end{enumerate}
